\documentclass[12pt,a4paper]{report}

\usepackage[spanish]{babel}
\usepackage{iftex}
\ifPDFTeX
\usepackage[utf8]{inputenc}
\usepackage[T1]{fontenc}
\usepackage[scaled=0.95]{helvet}
\renewcommand{\familydefault}{\sfdefault}
\else
\usepackage{fontspec}
\setmainfont{Arial}
\setsansfont{Arial}
\fi
\usepackage{geometry}
\geometry{margin=2.5cm}
\usepackage{setspace}
\setstretch{1.25}
\usepackage{booktabs}
\usepackage{longtable}
\usepackage{array}
\usepackage{enumitem}
\usepackage{hyperref}
\usepackage{xcolor}
\usepackage{titlesec}

\hypersetup{
	colorlinks=true,
	linkcolor=blue,
	urlcolor=blue,
	citecolor=blue,
	pdftitle={Especificación resumida de infraestructura y operación},
	pdfauthor={Luisdavid Colina; Susana Carvallo}
}

\titleformat{\section}
	{\bfseries\Large}
	{\thesection}
	{0.6em}
	{\Large}

\titleformat{\subsection}
	{\bfseries\normalsize}
	{\thesubsection}
	{0.6em}
	{\large}

\titlespacing*{\chapter}{0pt}{-8pt}{18pt}
\titlespacing*{\section}{0pt}{14pt}{8pt}
\titlespacing*{\subsection}{0pt}{10pt}{6pt}

\setlist[itemize]{itemsep=0.3em, topsep=0.3em}
\setlist[enumerate]{itemsep=0.3em, topsep=0.3em}

\pagestyle{plain}

\begin{document}

\begin{center}
{\LARGE\bfseries Especificación resumida}
\end{center}
\addcontentsline{toc}{chapter}{Especificación resumida}

\section*{Propósito operativo}
La solución debe poder ejecutarse con una base técnica simple: R 4.2.0 o superior, PostgreSQL 13+ como base de datos objetivo y Docker como opción.
En términos prácticos, esto significa que cada componente cumple una función concreta: el servidor ejecuta la aplicación, la base de datos guarda la información organizada, el almacenamiento conserva los archivos, la red permite acceso seguro y el monitoreo mantiene el servicio sano.
Para una puesta en marcha rápida y de bajo costo, la base de datos puede ir en la misma máquina que la aplicación al inicio; si luego crece el uso, se puede separar más adelante sin rehacer todo el sistema.

\subsection*{Especificaciones de Hardware de Referencia}

Para garantizar la integridad operativa y la eficiencia en la ejecución del sistema, se definen los siguientes parámetros de hardware inicial:

\begin{itemize}
    \item \textbf{Capacidad de Cómputo:} 1 a 2 vCPU como base técnica. Esta configuración asegura un procesamiento concurrente estable para las tareas de indexación y servicio de peticiones, manteniendo un perfil de consumo de recursos optimizado.
    \item \textbf{Memoria RAM:} Configuración base de 1 a 2 GB, con un escalado recomendado a 4 GB para entornos de producción muy demandantes. Este margen técnico permite la gestión eficiente de la caché de base de datos y la carga de objetos en el entorno R sin comprometer el rendimiento del sistema operativo.
    \item \textbf{Almacenamiento y Persistencia:} 40 GB en estado sólido (SSD) para el volumen de sistema. La arquitectura contempla la segregación de activos digitales en un volumen lógico independiente, permitiendo el crecimiento asíncrono del repositorio documental y mitigando riesgos de saturación en la partición raíz.
\end{itemize}
\subsection*{Servidor principal}
El servidor de aplicación es la computadora donde corre el sistema. Allí se realizan las consultas, los inicios de sesión, los filtros, la navegación y las exportaciones.
\begin{longtable}{p{0.28\textwidth} p{0.62\textwidth}}
\toprule
\textbf{Elemento} & \textbf{Especificación y explicación} \\
\midrule
CPU & 1 a 2 vCPU según el escenario. Como recomendación práctica, 2 vCPU da un margen cómodo para operación estable. \\
RAM & 1 GB para desarrollo mínimo y 2 GB como operación objetivo. Si se quiere más holgura, 4 GB es una cifra razonable para producción demandante. \\
Disco & 40 GB SSD como base operativa, ampliable según el crecimiento del repositorio. \\
\bottomrule
\end{longtable}

\subsection*{Almacenamiento de archivos}
El almacenamiento documental sirve para guardar los archivos digitales del repositorio aparte del sistema operativo. Así, si el programa se reinstala, los documentos siguen seguros y el sistema no se vuelve lento por falta de espacio.
El disco de almacenamiento puede ampliarse cuando crezca el repositorio.
También debe existir un volumen o espacio separado y exclusivo para los archivos escaneados, de modo que no se mezclen con otros datos del sistema.
\begin{longtable}{p{0.28\textwidth} p{0.62\textwidth}}
\toprule
\textbf{Elemento} & \textbf{Especificación y explicación} \\
\midrule
Volumen para archivos & Debe estar separado del sistema operativo y crecer de forma independiente. Esto evita que el aumento de PDFs comprometa la estabilidad del servidor. \\
Disco & 40 GB SSD como base operativa, con posibilidad de ampliar si el repositorio crece. \\
Estructura de respaldo & Debe existir un destino adicional para copias de seguridad. El repositorio no debe depender de una sola ubicación física. \\
Retención & Se recomienda conservar varias generaciones de respaldo para poder recuperar información ante errores humanos o corrupción de datos. \\
\bottomrule
\end{longtable}

\subsection*{Base de datos local}
La base de datos es donde se guardan los datos organizados: usuarios, nombres, fechas, categorías y relaciones entre documentos. No guarda la pantalla del sistema ni los PDF; guarda la información que permite buscar y ordenar todo.
Si se busca rapidez de despliegue y menor costo, se puede montar junto al sistema en la misma máquina.
\begin{longtable}{p{0.28\textwidth} p{0.62\textwidth}}
\toprule
\textbf{Elemento} & \textbf{Especificación y explicación} \\
\midrule
Motor & PostgreSQL 13 o superior, como base de datos objetivo. \\
RAM dedicada & Si la base de datos corre en la misma máquina, una reserva práctica de 2 GB adicionales es razonable para mantener el sistema fluido. Si se mantiene todo junto al inicio, conviene dejar margen para que no se vuelva lento. \\
Persistencia & La base debe sobrevivir reinicios y despliegues, por lo que no puede depender de almacenamiento temporal. \\
\bottomrule
\end{longtable}

\subsection*{Conectividad y acceso}
La red y el acceso controlan cómo entra la gente al sistema y cómo se administra. Su función es que el sistema pueda abrirse de forma segura, sin mostrar cosas que no deben estar visibles.
\begin{longtable}{p{0.28\textwidth} p{0.62\textwidth}}
\toprule
\textbf{Elemento} & \textbf{Especificación y explicación} \\
\midrule
Conectividad & Acceso estable a la red institucional. Si el sistema se publica externamente, debe hacerlo por HTTPS. \\
Firewall & Solo deben quedar expuestos los puertos necesarios para la aplicación y la administración. Reduce superficie de ataque y errores operativos. \\
Acceso remoto & La administración del servidor debe hacerse por canales restringidos, con credenciales individuales y registro de actividad. \\
\bottomrule
\end{longtable}

\section*{Mantenimiento operativo}
\subsection*{Copias de seguridad}
Debe existir una copia automática de los datos, la configuración y los archivos. Sirve para recuperar el repositorio si ocurre una falla, un error humano o una pérdida de información.

\subsection*{Seguimiento básico}
Debe haber una revisión básica del estado del sistema: si está encendido, si le queda espacio, si usa demasiada memoria o si está fallando. Eso permite actuar antes de que el servicio se caiga.

\subsection*{Actualización controlada}
El entorno debe permitir actualizaciones controladas del programa y sus componentes. Eso ayuda a corregir fallas y evita que el sistema se quede viejo o incompatible.

\subsection*{Guía técnica}
Se debe conservar una guía básica de instalación, configuración y recuperación. Esa guía permite que otra persona pueda operar o restaurar el sistema sin depender de que alguien se acuerde de todo.

\section*{Criterios de sostenibilidad}
\begin{itemize}
	\item La plataforma debe poder mantenerse con infraestructura de bajo costo. Así no se vuelve una carga difícil de pagar o sostener.
	\item La base de datos puede ir junto al sistema al inicio. Eso reduce costo y acelera el arranque.
	\item La arquitectura debe permitir crecer sin cambiar todo desde cero. Eso facilita sumar usuarios, documentos o módulos sin empezar de nuevo.
	\item La separación entre aplicación, datos y archivos debe ser clara. Eso ayuda a cuidar mejor la información y a arreglar cada parte por separado.
	\item La solución debe ser administrable por personal institucional con apoyo técnico limitado. Eso garantiza que el sistema siga funcionando aunque no haya un equipo grande de soporte.
\end{itemize}

\section*{Criterios de dimensionamiento}
Las cifras que no están fijadas de forma exacta se pusieron como estimaciones operativas realistas. Sirven para orientar la compra o el despliegue, pero pueden ajustarse si luego se define una carga de usuarios o un volumen documental más preciso.

\section*{Resumen operativo}
La prioridad no es construir un sistema grande, sino uno que pueda operar con recursos razonables, recuperarse ante fallos y mantenerse vigente con mantenimiento simple. Por eso los requisitos se enfocan en estabilidad, almacenamiento separado, respaldo, monitoreo y posibilidad de crecimiento gradual.

\end{document}
